% Contenido del Anexo 1 - Modelo de Casos de Uso

% Primera página: Información del proyecto y tabla de revisiones
\subsection{Proyecto Laboratorio de Semillas INIA}

\subsubsection{Modelo de Casos de Uso}

\vspace{1em}
\textbf{Versión [1.2]}

\vspace{2em}
\paragraph{Historia de revisiones}

\begin{table}[h]
\centering
\begin{tabular}{|p{3cm}|p{2cm}|p{6cm}|p{3cm}|}
\hline
\textbf{Fecha} & \textbf{Versión} & \textbf{Descripción} & \textbf{Autor} \\
\hline
10/09/2025 & 1.1 & Se definieron los casos de uso para el sistema a realizar & Nahuel Perdomo, Milagros Vairo, Nadia Gorría \\
\hline
17/11/2025 & 1.2 & Se redefinieron los casos de uso en base a las implementaciones finales & Nahuel Perdomo, Milagros Vairo, Nadia Gorría \\
\hline
\end{tabular}
\end{table}

\clearpage

% Segunda página: Actores
\subsection{Actores}
En esta sección se describe a cada uno de los actores que existen en el sistema. Un actor es un usuario del sistema o cualquier otro sistema que interactúa con el mismo.

\subsubsection{Analista}
\textbf{Descripción:} Usuario de laboratorio que registra muestras y carga los resultados de los análisis (Pureza, Germinación, PMS, Tetrazolio, etc.). Se encarga de ingresar datos en la plataforma y asegurar que las mediciones queden registradas correctamente.

\subsubsection{Administrador}
\textbf{Descripción:} Usuario con acceso completo al sistema. Además de las funciones del analista, puede gestionar usuarios, catálogos y normativas, y validar lotes. Es responsable de la administración general del sistema y de la generación de reportes oficiales.

\subsubsection{Observador}
\textbf{Descripción:} Usuario externo (cliente del laboratorio) o interno con acceso de solo lectura, que consulta los resultados de sus lotes. Puede descargar reportes y visualizar el estado de los análisis, pero no modificar información.

\clearpage

% Tercera página: Casos de Uso
\subsection{Casos de Uso}
En esta sección se describen los casos de uso identificados para el sistema. Un caso de uso representa una funcionalidad o interacción específica entre un actor y el sistema, detallando los pasos que se siguen para cumplir un objetivo particular.

Cada caso de uso se documenta con su descripción, precondiciones, flujo principal de eventos, flujos alternativos, postcondiciones y requerimientos especiales. Estos casos permiten comprender cómo los actores interactúan con el sistema y sirven de base para el diseño y validación del mismo.

\subsection{Caso de uso: Registrar Lote de Semillas}

\textbf{Descripción:} El usuario registra un nuevo lote de semillas con su información completa, asignándole los tipos de análisis que se realizarán.

\vspace{0.5em}
\noindent\textbf{PreCondiciones:}
\begin{itemize}[nosep]
    \item El usuario debe estar autenticado como Analista o Administrador.
    \item Debe existir al menos un cultivar en el sistema.
    \item Debe existir al menos una empresa registrada.
    \item Debe existir al menos un cliente registrado.
\end{itemize}

\vspace{0.5em}
\noindent\textbf{PostCondiciones:}
\begin{itemize}[nosep]
    \item Lote creado con estado activo = true.
    \item Lote disponible para crear análisis de los tipos asignados.
    \item Lote visible en el listado de lotes activos.
\end{itemize}

\vspace{0.5em}
\noindent\textbf{Actores:} Analista, Administrador.

\vspace{0.5em}
\noindent\textbf{Flujo de eventos principal:}
\begin{enumerate}[nosep]
    \item El usuario accede a la página ``Registro de Lotes''.
    \item Completa los datos obligatorios:
    \begin{itemize}[nosep]
        \item Ficha (campo único)
        \item Nombre del lote (campo único)
        \item Cultivar (selección desde catálogo)
        \item Tipo de lote (INTERNO/EXTERNO)
        \item Empresa (selección desde lista de contactos)
        \item Cliente (selección desde lista de contactos)
        \item Fecha de recibo
        \item Kilos limpios
    \end{itemize}
    \item Completa datos opcionales:
    \begin{itemize}[nosep]
        \item Código CC, Código FF
        \item Fecha de entrega
        \item Depósito, Unidad de embolsado, Remitente
        \item Datos de humedad (tipo y valor)
        \item Número de artículo, Origen, Estado
        \item Fecha de cosecha
        \item Observaciones
    \end{itemize}
    \item Selecciona los tipos de análisis que se realizarán al lote:
    \begin{itemize}[nosep]
        \item Pureza
        \item PMS (Peso de Mil Semillas)
        \item Germinación
        \item Tetrazolio
        \item DOSN (Determinación de Otras Semillas en Número)
    \end{itemize}
    \item El sistema valida que:
    \begin{itemize}[nosep]
        \item La ficha no esté duplicada
        \item El nombre del lote no esté duplicado
        \item La fecha de recibo no sea futura
        \item La fecha de cosecha no sea futura
    \end{itemize}
    \item El sistema crea el lote con \texttt{activo = true}.
    \item El lote aparece en la lista de ``Últimos lotes registrados''.
    \item El sistema muestra mensaje de confirmación.
\end{enumerate}

\vspace{0.5em}
\noindent\textbf{Flujos alternativos:}

\vspace{0.3em}
\noindent\textit{1. Ficha duplicada}
\begin{itemize}[nosep]
    \item El sistema valida en tiempo real
    \item Muestra alerta: ``Esta ficha ya está registrada. Por favor, utiliza una ficha diferente''
    \item No permite guardar hasta corregir
\end{itemize}

\vspace{0.3em}
\noindent\textit{2. Nombre de lote duplicado}
\begin{itemize}[nosep]
    \item El sistema valida en tiempo real
    \item Muestra alerta: ``Este nombre de lote ya está registrado. Por favor, utiliza un nombre diferente''
    \item No permite guardar hasta corregir
\end{itemize}

\vspace{0.3em}
\noindent\textit{3. Fecha futura}
\begin{itemize}[nosep]
    \item El sistema valida en backend
    \item Muestra error: ``La fecha de recibo/cosecha no puede ser posterior a la fecha actual''
    \item No permite guardar hasta corregir
\end{itemize}

\vspace{0.3em}
\noindent\textit{4. Sin tipos de análisis asignados}
\begin{itemize}[nosep]
    \item El lote se crea sin restricciones
    \item No se podrán crear análisis hasta editar el lote y asignar tipos
\end{itemize}

\vspace{0.5em}
\noindent\textbf{Requerimientos especiales:}
\begin{itemize}[nosep]
    \item Validación asíncrona de campos únicos (ficha y nombre)
    \item Eliminación automática de duplicados en tipos de análisis asignados
    \item Los tipos de análisis asignados se almacenan como enum JSON
\end{itemize}

\subsubsection{Registro de Análisis Pureza}

\noindent\textbf{Caso de uso:} Cargar Resultados de Pureza

\vspace{0.5em}
\noindent\textbf{Descripción:} El analista registra pesos y componentes de una muestra de semillas para calcular el porcentaje de pureza física, identificando semilla pura, materia inerte, otros cultivos y malezas.

\vspace{0.5em}
\noindent\textbf{PreCondición:}
\begin{itemize}[nosep]
    \item Debe existir al menos un lote activo con tipo de análisis ``PUREZA'' asignado
    \item El lote no debe tener un análisis de Pureza en estado APROBADO
    \item El usuario debe estar autenticado como Analista o Administrador
\end{itemize}

\vspace{0.5em}
\noindent\textbf{PostCondición:}
\begin{itemize}[nosep]
    \item Análisis creado con estado \texttt{EN\_PROCESO}
    \item Fecha de inicio automática asignada
    \item Resultados guardados y vinculados al lote
    \item Registro automático en historial de análisis
\end{itemize}

\vspace{0.5em}
\noindent\textbf{Actores:} Analista, Administrador.

\vspace{0.5em}
\noindent\textbf{Flujo de evento principal:}
\begin{enumerate}[nosep]
    \item \textbf{Fase 1: Creación del análisis}
    \begin{itemize}[nosep]
        \item El usuario selecciona un lote elegible para Pureza
        \item Accede a ``Crear análisis de Pureza''
        \item El sistema crea el análisis con estado \texttt{EN\_PROCESO} y fecha de inicio automática
    \end{itemize}
    \item \textbf{Fase 2: Registro de datos INIA}
    \begin{itemize}[nosep]
        \item Ingresa fecha del análisis (obligatoria)
        \item Ingresa peso inicial en gramos (obligatorio, debe ser > 0)
        \item Ingresa componentes en gramos: Semilla pura, Materia inerte, Otros cultivos, Malezas, Malezas toleradas, Malezas tolerancia cero
        \item El sistema calcula automáticamente: peso total y porcentajes sin redondeo (4 decimales)
        \item Validaciones en tiempo real: peso total $\leq$ peso inicial; alerta si pérdida de masa > 5\% (informativa)
        \item El usuario ingresa porcentajes redondeados manualmente; el sistema valida suma = 100\%
        \item Valores entre 0 y 0.05\% muestran ``tr'' (trazas) y son de solo lectura
    \end{itemize}
    \item \textbf{Fase 3: Datos INASE (opcionales)}
    \begin{itemize}[nosep]
        \item El usuario puede ingresar fecha INASE y porcentajes oficiales INASE
    \end{itemize}
    \item \textbf{Fase 4: Registro de listados (opcional)}
    \begin{itemize}[nosep]
        \item Registro de malezas identificadas y otros cultivos, indicando catálogo/especie, cantidad y categorías
    \end{itemize}
    \item \textbf{Fase 5: Cumplimiento del estándar}
    \begin{itemize}[nosep]
        \item El usuario selecciona si cumple o no con el estándar y el sistema guarda los cambios
    \end{itemize}
\end{enumerate}

\vspace{0.5em}
\noindent\textbf{Flujos alternativos:}
\begin{itemize}[nosep]
    \item \textit{Peso total mayor que peso inicial}: bloqueo de guardado y mensaje de error
    \item \textit{Pérdida de masa excesiva}: muestra alerta informativa (no bloquea)
    \item \textit{Suma de porcentajes redondeados $\neq$ 100\%}: muestra alerta en rojo (no bloquea)
    \item \textit{Finalizar sin evidencia}: sistema rechaza la finalización si faltan datos
\end{itemize}

\vspace{0.5em}
\noindent\textbf{Requerimientos especiales:}
\begin{itemize}[nosep]
    \item Cálculos automáticos en tiempo real con 4 decimales
    \item Validación de coherencia en pesos (total $\leq$ inicial)
    \item Alerta de pérdida de masa > 5\% y manejo de ``tr'' para trazas
    \item Tabla de tolerancias disponible para consulta
\end{itemize}

\subsubsection{Registro de Análisis PMS}

\noindent\textbf{Caso de uso:} Configurar Análisis de Peso de Mil Semillas

\vspace{0.5em}
\noindent\textbf{Descripción:} El usuario configura los parámetros iniciales para un análisis de PMS (Peso de Mil Semillas), estableciendo el número de repeticiones esperadas y el tipo de semilla.

\vspace{0.5em}
\noindent\textbf{PreCondición:}
\begin{itemize}[nosep]
    \item Debe existir al menos un lote activo con tipo de análisis ``PMS'' asignado
    \item El lote no debe tener un análisis de PMS en estado \texttt{APROBADO}
    \item El usuario debe estar autenticado como Analista o Administrador
\end{itemize}

\vspace{0.5em}
\noindent\textbf{PostCondición:}
\begin{itemize}[nosep]
    \item Análisis creado con estado \texttt{EN\_PROCESO}
    \item Configuración de repeticiones establecida
    \item Sistema listo para recibir tandas y repeticiones
    \item Umbral de coeficiente de variación establecido según tipo de semilla
\end{itemize}

\vspace{0.5em}
\noindent\textbf{Actores:} Analista, Administrador.

\vspace{0.5em}
\noindent\textbf{Flujo de evento principal:}
\begin{enumerate}[nosep]
    \item El usuario selecciona un lote elegible para PMS
    \item Accede a ``Crear análisis de PMS''
    \item Configura parámetros obligatorios:
    \begin{itemize}[nosep]
        \item Número de repeticiones esperadas: entre 4 y 20 (obligatorio)
        \item Tipo de semilla: marca checkbox si es semilla brozosa (opcional)
    \end{itemize}
    \item El sistema muestra:
    \begin{itemize}[nosep]
        \item Número de tandas inicial = 1 (fijo, no editable)
        \item Umbral CV: $\leq$ 6\% si es brozosa, $\leq$ 4\% si es normal
    \end{itemize}
    \item El usuario guarda la configuración
    \item El sistema:
    \begin{itemize}[nosep]
        \item Crea el análisis con estado \texttt{EN\_PROCESO}
        \item Crea automáticamente la primera tanda vacía
        \item Queda listo para agregar repeticiones
    \end{itemize}
\end{enumerate}

\vspace{0.5em}
\noindent\textbf{Flujos alternativos:}
\begin{itemize}[nosep]
    \item \textit{Número de repeticiones inválido:}
    \begin{itemize}[nosep]
        \item El sistema valida que esté entre 4 y 20
        \item Muestra error: ``Debe especificar un número válido de repeticiones (4-20)''
        \item No permite guardar hasta corregir
    \end{itemize}
    \item \textit{Lote sin tipo de análisis PMS asignado:}
    \begin{itemize}[nosep]
        \item El sistema no permite seleccionar el lote
        \item Solo muestra lotes elegibles con PMS asignado
    \end{itemize}
\end{itemize}

\vspace{0.5em}
\noindent\textbf{Requerimientos especiales:}
\begin{itemize}[nosep]
    \item Validación estricta de rango de repeticiones (4-20)
    \item Umbral de CV dinámico según tipo de semilla (4\% o 6\%)
    \item Número de repeticiones no modificable después de crear análisis
    \item Primera tanda se crea automáticamente al registrar el análisis
\end{itemize}

\vspace{0.5em}
\noindent\textbf{Nota:} La creación de repeticiones, carga de datos y cálculos (peso promedio, CV, creación automática de tandas adicionales) se realiza en la fase de edición del análisis. El registro solo configura los parámetros iniciales.

\subsection{Caso de uso: Registro de Análisis de Germinación}

\noindent\textbf{Caso de uso:} Configurar Análisis de Germinación

\vspace{0.5em}
\noindent\textbf{Descripción:} El usuario configura un análisis de germinación, el cual servirá como contenedor para múltiples tablas de germinación con condiciones experimentales específicas.

\vspace{0.5em}
\noindent\textbf{PreCondición:}
\begin{itemize}[nosep]
    \item Debe existir al menos un lote activo con tipo de análisis ``GERMINACION'' asignado
    \item El lote no debe tener un análisis de Germinación en estado \texttt{APROBADO}
    \item El usuario debe estar autenticado como Analista o Administrador
\end{itemize}

\vspace{0.5em}
\noindent\textbf{PostCondición:}
\begin{itemize}[nosep]
    \item Análisis creado con estado \texttt{EN\_PROCESO}
    \item Sistema listo para crear tablas de germinación individuales
    \item Cada tabla tendrá su propia configuración de fechas, repeticiones y condiciones
\end{itemize}

\vspace{0.5em}
\noindent\textbf{Actores:} Analista, Administrador.

\vspace{0.5em}
\noindent\textbf{Flujo de evento principal:}
\begin{enumerate}[nosep]
    \item El usuario selecciona un lote elegible para Germinación
    \item Accede a ``Crear análisis de Germinación''
    \item El sistema muestra mensaje informativo:
    \begin{itemize}[nosep]
        \item ``Solo necesitas seleccionar el lote y opcionalmente agregar comentarios''
        \item ``La configuración detallada se realiza al crear cada tabla''
    \end{itemize}
    \item El usuario puede agregar comentarios generales (opcional)
    \item Guarda el análisis
    \item El sistema:
    \begin{itemize}[nosep]
        \item Crea el análisis con estado \texttt{EN\_PROCESO}
        \item Queda listo para crear tablas de germinación
    \end{itemize}
\end{enumerate}

\vspace{0.5em}
\noindent\textbf{Flujos alternativos:}
\begin{itemize}[nosep]
    \item \textit{Lote sin tipo de análisis GERMINACION asignado:}
    \begin{itemize}[nosep]
        \item El sistema no permite seleccionar el lote
        \item Solo muestra lotes elegibles con GERMINACION asignado
    \end{itemize}
    \item \textit{Guardar sin comentarios:}
    \begin{itemize}[nosep]
        \item El análisis se crea exitosamente
        \item Los comentarios son completamente opcionales
    \end{itemize}
\end{itemize}

\vspace{0.5em}
\noindent\textbf{Requerimientos especiales:}
\begin{itemize}[nosep]
    \item Análisis de germinación funciona como contenedor de tablas
    \item No requiere configuración inicial compleja
    \item Mínima configuración en el registro
    \item Tabla de tolerancias disponible para consulta
\end{itemize}

\vspace{0.5em}
\noindent\textbf{Nota:} La creación de tablas de germinación, configuración de condiciones (fechas, repeticiones, semillas, prefrío, pretratamiento, método, temperatura), creación de repeticiones, registro de conteos y cálculos automáticos se realiza en la fase de edición del análisis. El registro solo crea el contenedor inicial.

\subsubsection{Caso de uso: Registro de Análisis de Tetrazolio}

\noindent\textbf{Caso de uso:} Configurar Análisis de Tetrazolio

\vspace{0.5em}
\noindent\textbf{Descripción:} El usuario configura los parámetros del ensayo de viabilidad con tetrazolio, incluyendo pretratamiento, condiciones de tinción y número de repeticiones.

\vspace{0.5em}
\noindent\textbf{PreCondición:}
\begin{itemize}[nosep]
    \item Debe existir al menos un lote activo con tipo de análisis ``TETRAZOLIO'' asignado.
    \item El lote no debe tener un análisis de Tetrazolio en estado \texttt{APROBADO}.
    \item El usuario debe estar autenticado como Analista o Administrador.
\end{itemize}

\vspace{0.5em}
\noindent\textbf{PostCondición:}
\begin{itemize}[nosep]
    \item Análisis creado con estado \texttt{EN\_PROCESO}.
    \item Parámetros del ensayo configurados.
    \item Sistema listo para recibir resultados de repeticiones.
    \item Configuración no modificable después de crear repeticiones.
\end{itemize}

\vspace{0.5em}
\noindent\textbf{Actores:} Analista, Administrador.

\vspace{0.5em}
\noindent\textbf{Flujo de evento principal:}
\begin{enumerate}[nosep]
    \item \textbf{Configuración del ensayo}
    \begin{itemize}[nosep]
        \item El usuario selecciona un lote elegible para Tetrazolio
        \item Accede a ``Crear análisis de Tetrazolio''
        \item Configura parámetros obligatorios:
        \begin{itemize}[nosep]
            \item Fecha del ensayo
            \item Semillas por repetición: 25, 50 o 100 (no modificable después)
            \item Número de repeticiones esperadas: entre 2 y 8 (no modificable después)
            \item Pretratamiento: EP 16h, EP 18h, S/Pretratamiento, Agua 7h, Agua 8h, u Otro
            \item Concentración: 1\%, 0\%, 5\%, 0.75\%, u Otro
            \item Temperatura de tinción: 30-40°C u Otra
            \item Tiempo de tinción: 2h, 3h, 16h, 18h u Otra
        \end{itemize}
        \item Si selecciona ``Otro'' en algún campo, debe especificar el valor manualmente
        \item Ingresa Viabilidad INASE (opcional, numérico)
        \item Guarda la configuración
        \item El sistema crea el análisis con estado \texttt{EN\_PROCESO}
    \end{itemize}
\end{enumerate}

\vspace{0.5em}
\noindent\textbf{Flujos alternativos:}
\begin{itemize}[nosep]
    \item \textit{Campos ``Otro'' sin especificar:}
    \begin{itemize}[nosep]
        \item El sistema valida que si se selecciona ``Otro (especificar)'', el campo de texto no esté vacío
        \item Muestra error: ``Debe especificar el pretratamiento/concentración/temperatura/horas''
        \item No permite guardar hasta completar
    \end{itemize}
    \item \textit{Viabilidad INASE inválida:}
    \begin{itemize}[nosep]
        \item El sistema valida que sea un número válido
        \item Muestra error: ``El campo 'Viabilidad INASE' debe ser un número válido''
        \item No bloquea guardado si no se ingresa (es opcional)
    \end{itemize}
    \item \textit{Lote sin tipo de análisis TETRAZOLIO asignado:}
    \begin{itemize}[nosep]
        \item El sistema no permite seleccionar el lote
        \item Solo muestra lotes elegibles con TETRAZOLIO asignado
    \end{itemize}
\end{itemize}

\vspace{0.5em}
\noindent\textbf{Requerimientos especiales:}
\begin{itemize}[nosep]
    \item Validación estricta de rangos (repeticiones 2-8, semillas 25/50/100)
    \item Campos de configuración no modificables después de crear repeticiones
    \item Validación de campos ``Otro'' con texto obligatorio
    \item Tabla de tolerancias disponible para consulta
\end{itemize}

\vspace{0.5em}
\noindent\textbf{Nota:} La creación de repeticiones, ingreso de conteos (semillas viables, no viables, muertas) y cálculos automáticos (totales, porcentajes de viabilidad, promedios) se realiza en la fase de edición del análisis. El registro solo configura los parámetros del ensayo.

\subsubsection{Registro de Análisis DOSN}

\noindent\textbf{Caso de uso:} Registrar Determinación de Otras Semillas en Número

\vspace{0.5em}
\noindent\textbf{Descripción:} El analista registra los resultados del análisis DOSN (Determinación de Otras Semillas en Número) para INIA y opcionalmente para INASE, identificando malezas, otros cultivos y semillas brassica encontradas.

\vspace{0.5em}
\noindent\textbf{PreCondición:}
\begin{itemize}[nosep]
    \item Debe existir al menos un lote activo con tipo de análisis ``DOSN'' asignado.
    \item El lote no debe tener un análisis DOSN en estado \texttt{APROBADO}.
    \item El usuario debe estar autenticado como Analista o Administrador.
\end{itemize}

\vspace{0.5em}
\noindent\textbf{PostCondición:}
\begin{itemize}[nosep]
    \item Análisis creado con estado \texttt{EN\_PROCESO}.
    \item Datos INIA registrados (obligatorios).
    \item Datos INASE registrados (opcionales).
    \item Listados de malezas, cultivos y brassica vinculados.
\end{itemize}

\vspace{0.5em}
\noindent\textbf{Actores:} Analista, Administrador.

\vspace{0.5em}
\noindent\textbf{Flujo de evento principal:}
\begin{enumerate}[nosep]
    \item \textbf{Configuración INIA (obligatoria)}
    \begin{itemize}[nosep]
        \item El usuario selecciona un lote elegible para DOSN.
        \item Accede a ``Crear análisis de DOSN''.
        \item Completa datos INIA obligatorios:
        \begin{itemize}[nosep]
            \item Fecha de análisis (no futura).
            \item Gramos analizados (> 0).
            \item Tipo de análisis: Completo, Reducido, Limitado o ReducidoLimitado (al menos uno).
        \end{itemize}
        \item El sistema valida en tiempo real.
    \end{itemize}
    \item \textbf{Configuración INASE (opcional)}
    \begin{itemize}[nosep]
        \item El usuario puede completar datos INASE:
        \begin{itemize}[nosep]
            \item Fecha de análisis.
            \item Gramos analizados.
            \item Tipos de análisis (mismos que INIA).
        \end{itemize}
        \item Si no completa, el análisis se crea solo con datos INIA.
    \end{itemize}
    \item \textbf{Registro de listados (opcional)}
    \begin{itemize}[nosep]
        \item El usuario puede registrar:
        \begin{itemize}[nosep]
            \item Malezas (tolerancia cero, toleradas o comunes).
            \item Otros cultivos.
            \item Brassica (si aplica).
        \end{itemize}
        \item Para cada registro ingresa:
        \begin{itemize}[nosep]
            \item Catálogo o especie.
            \item Cantidad de semillas.
            \item Categoría aplicable.
        \end{itemize}
    \end{itemize}
    \item \textbf{Cuscuta (opcional)}
    \begin{itemize}[nosep]
        \item Si encuentra semillas de Cuscuta, registra:
        \begin{itemize}[nosep]
            \item Cantidad de semillas.
            \item Observaciones.
        \end{itemize}
    \end{itemize}
    \item \textbf{Cumplimiento del estándar}
    \begin{itemize}[nosep]
        \item El usuario selecciona si cumple o no con el estándar.
        \item Guarda el análisis.
    \end{itemize}
\end{enumerate}

\vspace{0.5em}
\noindent\textbf{Flujos alternativos:}
\begin{itemize}[nosep]
    \item \textit{Fecha futura:}
    \begin{itemize}[nosep]
        \item El sistema valida que la fecha INIA no sea posterior a hoy.
        \item Muestra error: ``Ingrese una fecha válida (no futura)''.
        \item No permite guardar hasta corregir.
    \end{itemize}
    \item \textit{Gramos inválidos:}
    \begin{itemize}[nosep]
        \item El sistema valida que los gramos sean > 0.
        \item Muestra error: ``Debe ingresar una cantidad válida de gramos (> 0)''.
        \item No permite guardar hasta corregir.
    \end{itemize}
    \item \textit{Sin tipo de análisis INIA:}
    \begin{itemize}[nosep]
        \item El sistema valida que al menos un checkbox esté marcado.
        \item Muestra error: ``Debe seleccionar al menos un tipo de análisis''.
        \item No permite guardar hasta corregir.
    \end{itemize}
    \item \textit{Finalizar sin evidencia:}
    \begin{itemize}[nosep]
        \item Si intenta finalizar sin listados, cuscuta ni datos suficientes.
        \item El sistema valida que haya algún tipo de evidencia registrada.
    \end{itemize}
\end{itemize}

\vspace{0.5em}
\noindent\textbf{Requerimientos especiales:}
\begin{itemize}[nosep]
    \item Validación de fechas no futuras.
    \item Validación de gramos > 0.
    \item Al menos un tipo de análisis INIA obligatorio.
    \item Datos INASE completamente opcionales.
    \item Tabla de tolerancias disponible para consulta.
    \item Permite múltiples categorías de listados.
\end{itemize}

\subsubsection{Caso de uso: Registrar Usuario}

\noindent\textbf{Caso de uso:} Registrar Solicitud de Usuario

\vspace{0.5em}
\noindent\textbf{Descripción:} Una persona solicita acceso al sistema completando un formulario de registro. La solicitud queda pendiente hasta ser aprobada por un administrador.

\vspace{0.5em}
\noindent\textbf{PreCondición:}
\begin{itemize}[nosep]
    \item El solicitante no debe tener cuenta en el sistema
    \item El nombre de usuario y email deben ser únicos
\end{itemize}

\vspace{0.5em}
\noindent\textbf{PostCondición:}
\begin{itemize}[nosep]
    \item Usuario creado con estado \texttt{PENDIENTE}
    \item Usuario con \texttt{activo = false} y \texttt{rol = null}
    \item Notificación automática creada para analistas
    \item Email de confirmación enviado al solicitante
    \item Email de notificación enviado a todos los analistas
\end{itemize}

\vspace{0.5em}
\noindent\textbf{Actores:} Solicitante (persona sin cuenta), Administrador (recibe notificación).

\vspace{0.5em}
\noindent\textbf{Flujo de evento principal:}
\begin{enumerate}[nosep]
    \item El solicitante accede a la página de registro desde el login
    \item Completa el formulario con datos obligatorios:
    \begin{itemize}[nosep]
        \item Nombre de usuario (único, mínimo 3 caracteres, solo letras, números, puntos y guiones bajos)
        \item Nombres (solo letras)
        \item Apellidos (solo letras)
        \item Email (único, formato válido)
        \item Contraseña (mínimo 8 caracteres, al menos una letra y un número)
        \item Confirmar contraseña (debe coincidir)
    \end{itemize}
    \item El sistema valida en tiempo real:
    \begin{itemize}[nosep]
        \item Formato de cada campo
        \item Unicidad de nombre de usuario
        \item Unicidad de email
        \item Fortaleza de contraseña (débil/media/fuerte)
        \item Coincidencia de contraseñas
    \end{itemize}
    \item Al perder el foco (blur), valida unicidad de nombre y email
    \item El solicitante envía el formulario
    \item El sistema:
    \begin{itemize}[nosep]
        \item Hashea la contraseña con bcrypt
        \item Crea el usuario con estado \texttt{PENDIENTE}, \texttt{activo = false}, \texttt{rol = null}
        \item Crea notificación automática para analistas
        \item Envía email de confirmación al solicitante
        \item Envía email a todos los analistas notificando la nueva solicitud
    \end{itemize}
    \item Muestra diálogo de éxito explicando el proceso de aprobación:
    \begin{itemize}[nosep]
        \item ``Esperando aprobación''
        \item ``El administrador debe aprobar tu solicitud antes de que puedas acceder''
        \item ``Recibirás una notificación por correo cuando sea aprobada''
    \end{itemize}
    \item Redirige al login al cerrar el diálogo
\end{enumerate}

\vspace{0.5em}
\noindent\textbf{Flujos alternativos:}
\begin{itemize}[nosep]
    \item \textit{Nombre de usuario duplicado:}
    \begin{itemize}[nosep]
        \item El sistema detecta en blur o en submit
        \item Muestra error: ``Este nombre de usuario ya está registrado''
        \item No permite guardar hasta cambiar
    \end{itemize}
    \item \textit{Email duplicado:}
    \begin{itemize}[nosep]
        \item El sistema detecta en blur o en submit
        \item Muestra error: ``Este correo electrónico ya está registrado''
        \item No permite guardar hasta cambiar
    \end{itemize}
    \item \textit{Contraseña débil:}
    \begin{itemize}[nosep]
        \item El sistema muestra indicador de fortaleza:
        \begin{itemize}[nosep]
            \item Rojo (débil): < 8 caracteres o sin letra/número
            \item Amarillo (media): 8+ caracteres con letra y número
            \item Verde (fuerte): 8+ caracteres con letra, número y caracteres especiales
        \end{itemize}
        \item Permite guardar con contraseña media o fuerte
    \end{itemize}
    \item \textit{Contraseñas no coinciden:}
    \begin{itemize}[nosep]
        \item El sistema valida en tiempo real
        \item Muestra error: ``Las contraseñas no coinciden''
        \item No permite guardar hasta corregir
    \end{itemize}
    \item \textit{Error al enviar emails:}
    \begin{itemize}[nosep]
        \item El registro se completa exitosamente
        \item Los emails se intentan enviar pero no bloquean el registro
        \item Se registra log de error pero no afecta al usuario
    \end{itemize}
\end{itemize}

\vspace{0.5em}
\noindent\textbf{Requerimientos especiales:}
\begin{itemize}[nosep]
    \item Validación asíncrona de unicidad (nombre y email)
    \item Validación en tiempo real de fortaleza de contraseña
    \item Hasheo seguro con bcrypt antes de almacenar
    \item Envío de emails asíncrono (no bloquea registro)
    \item Creación automática de notificaciones
    \item Indicador visual de fortaleza de contraseña (débil/media/fuerte)
\end{itemize}

\subsubsection{Edición de Lotes}

\noindent\textbf{Caso de uso:} Editar Lote de Semillas

\vspace{0.5em}
\noindent\textbf{Descripción:} El usuario modifica la información de un lote existente, pudiendo actualizar datos generales y tipos de análisis asignados, respetando las restricciones de análisis ya creados.

\vspace{0.5em}
\noindent\textbf{Pre-Condición:}
\begin{itemize}[nosep]
    \item El lote debe existir y estar activo
    \item El usuario debe estar autenticado como Analista o Administrador
    \item No se puede editar un lote inactivo sin reactivarlo primero
\end{itemize}

\vspace{0.5em}
\noindent\textbf{Post-Condición:}
\begin{itemize}[nosep]
    \item Información del lote actualizada
    \item Cambios reflejados inmediatamente en el sistema
    \item Si se agregaron nuevos tipos de análisis, quedan disponibles para crear
\end{itemize}

\vspace{0.5em}
\noindent\textbf{Actores:} Analista, Administrador.

\vspace{0.5em}
\noindent\textbf{Flujo de evento principal:}
\begin{enumerate}[nosep]
    \item El usuario accede al detalle de un lote activo
    \item Selecciona ``Editar lote''
    \item El sistema carga los datos actuales del lote
    \item El usuario modifica campos deseados:
    \begin{itemize}[nosep]
        \item Ficha (único)
        \item Nombre del lote (único)
        \item Cultivar
        \item Tipo de lote (INTERNO/EXTERNO)
        \item Empresa, Cliente
        \item Fechas (recibo, entrega, cosecha)
        \item Kilos limpios
        \item Datos de humedad
        \item Depósito, Unidad de embolsado, Remitente
        \item Número de artículo, Origen, Estado
        \item Observaciones
    \end{itemize}
    \item Puede agregar nuevos tipos de análisis
    \item Intenta remover tipos de análisis existentes
    \begin{itemize}[nosep]
        \item El sistema valida:
        \begin{itemize}[nosep]
            \item Si hay análisis creados de ese tipo $\rightarrow$ no permite remover
            \item Si no hay análisis $\rightarrow$ permite remover
        \end{itemize}
    \end{itemize}
    \item Valida unicidad de ficha y nombre de lote (si cambiaron)
    \item Valida fechas no futuras
    \item Guarda los cambios
    \item El sistema muestra confirmación
\end{enumerate}

\vspace{0.5em}
\noindent\textbf{Flujos alternativos:}
\begin{itemize}[nosep]
    \item \textit{Intentar editar lote inactivo:}
    \begin{itemize}[nosep]
        \item El sistema bloquea la edición
        \item Muestra error: ``No se puede editar un lote inactivo. Debe reactivarlo primero''
        \item Ofrece botón ``Reactivar lote''
    \end{itemize}
    \item \textit{Intentar remover tipo de análisis con análisis creados:}
    \begin{itemize}[nosep]
        \item El sistema detecta análisis existentes de ese tipo
        \item Muestra error: ``No se puede remover el tipo de análisis [TIPO] porque ya existen análisis creados de este tipo para el lote''
        \item El checkbox permanece marcado y deshabilitado
    \end{itemize}
    \item \textit{Ficha o nombre duplicado:}
    \begin{itemize}[nosep]
        \item El sistema valida en tiempo real y en submit
        \item Muestra alerta correspondiente
        \item No permite guardar hasta corregir
    \end{itemize}
    \item \textit{Fecha futura:}
    \begin{itemize}[nosep]
        \item El sistema valida en backend
        \item Muestra error: ``La fecha de recibo/cosecha no puede ser posterior a la fecha actual''
        \item No permite guardar hasta corregir
    \end{itemize}
\end{itemize}

\vspace{0.5em}
\noindent\textbf{Requerimientos especiales:}
\begin{itemize}[nosep]
    \item No se puede editar lote inactivo directamente
    \item Tipos de análisis solo se pueden agregar, no remover si hay análisis creados
    \item Validación asíncrona de campos únicos (ficha y nombre) excluyendo el lote actual
    \item Eliminación automática de duplicados en tipos de análisis asignados
\end{itemize}

\subsubsection{Gestión de Análisis de Pureza}

\noindent\textbf{Caso de uso:} Editar y Completar Análisis de Pureza

\vspace{0.5em}
\noindent\textbf{Descripción:} El usuario edita un análisis de Pureza existente, completa datos de pesos, porcentajes, registra listados de malezas y otros cultivos, y define cumplimiento del estándar.

\vspace{0.5em}
\noindent\textbf{Pre-Condición:}
\begin{itemize}[nosep]
    \item El análisis de Pureza debe existir.
    \item El análisis debe estar en estado \texttt{EN\_PROCESO} o \texttt{PENDIENTE\_APROBACION}.
    \item Si está \texttt{APROBADO}, solo Administrador puede editar.
    \item El usuario debe estar autenticado.
\end{itemize}

\vspace{0.5em}
\noindent\textbf{Post-Condición:}
\begin{itemize}[nosep]
    \item Datos del análisis actualizados.
    \item Si Analista edita \texttt{APROBADO} $\rightarrow$ cambia a \texttt{PENDIENTE\_APROBACION}.
    \item Si Administrador edita \texttt{APROBADO} $\rightarrow$ permanece \texttt{APROBADO}.
    \item Registro automático en historial de modificación.
\end{itemize}

\vspace{0.5em}
\noindent\textbf{Actores:} Analista, Administrador.

\vspace{0.5em}
\noindent\textbf{Flujo de evento principal:}
\begin{enumerate}[nosep]
    \item \textbf{Fase 1: Acceso al análisis}
    \begin{itemize}[nosep]
        \item El usuario accede al detalle del análisis de Pureza.
        \item Selecciona ``Editar análisis''.
        \item El sistema carga los datos actuales.
    \end{itemize}
    \item \textbf{Fase 2: Edición de datos INIA}
    \begin{itemize}[nosep]
        \item El usuario puede modificar:
        \begin{itemize}[nosep]
            \item Fecha del análisis.
            \item Peso inicial (g).
            \item Componentes en gramos (semilla pura, materia inerte, otros cultivos, malezas, malezas toleradas, malezas tol. cero).
        \end{itemize}
        \item El sistema recalcula automáticamente:
        \begin{itemize}[nosep]
            \item Peso total.
            \item Porcentajes sin redondeo (4 decimales).
        \end{itemize}
        \item El sistema valida en tiempo real:
        \begin{itemize}[nosep]
            \item Peso total $\leq$ peso inicial.
            \item Pérdida de masa > 5\% $\rightarrow$ alerta.
        \end{itemize}
        \item El usuario ajusta porcentajes redondeados.
        \item El sistema valida suma = 100\%.
    \end{itemize}
    \item \textbf{Fase 3: Edición de datos INASE}
    \begin{itemize}[nosep]
        \item El usuario puede modificar:
        \begin{itemize}[nosep]
            \item Fecha INASE.
            \item Porcentajes oficiales INASE.
        \end{itemize}
    \end{itemize}
    \item \textbf{Fase 4: Gestión de listados}
    \begin{itemize}[nosep]
        \item El usuario puede:
        \begin{itemize}[nosep]
            \item Agregar nuevas malezas identificadas.
            \item Agregar nuevos otros cultivos.
            \item Editar registros existentes.
            \item Eliminar registros.
        \end{itemize}
    \end{itemize}
    \item \textbf{Fase 5: Cumplimiento del estándar}
    \begin{itemize}[nosep]
        \item El usuario selecciona cumplimiento (Sí/No).
        \item Guarda los cambios.
        \item El sistema:
        \begin{itemize}[nosep]
            \item Actualiza el análisis.
            \item Registra modificación en historial.
            \item Si es Analista editando \texttt{APROBADO} $\rightarrow$ cambia a \texttt{PENDIENTE\_APROBACION}.
            \item Si es Administrador $\rightarrow$ mantiene estado actual.
        \end{itemize}
    \end{itemize}
\end{enumerate}

\vspace{0.5em}
\noindent\textbf{Flujos alternativos:}
\begin{itemize}[nosep]
    \item \textit{Analista edita análisis \texttt{APROBADO}:}
    \begin{itemize}[nosep]
        \item El sistema detecta que es Analista.
        \item Permite editar.
        \item Al guardar, cambia estado a \texttt{PENDIENTE\_APROBACION}.
        \item Muestra mensaje: ``El análisis ha sido devuelto a estado \texttt{PENDIENTE\_APROBACION} para nueva revisión''.
    \end{itemize}
    \item \textit{Administrador edita análisis \texttt{APROBADO}:}
    \begin{itemize}[nosep]
        \item El sistema detecta que es Administrador.
        \item Permite editar.
        \item Al guardar, mantiene estado \texttt{APROBADO}.
        \item Muestra mensaje: ``Análisis actualizado exitosamente''.
    \end{itemize}
    \item \textit{Peso total mayor que peso inicial:}
    \begin{itemize}[nosep]
        \item El sistema bloquea el guardado.
        \item Muestra error: ``El peso total no puede ser mayor al peso inicial''.
        \item El usuario debe corregir.
    \end{itemize}
    \item \textit{Suma de porcentajes $\neq$ 100\%:}
    \begin{itemize}[nosep]
        \item El sistema muestra alerta en rojo.
        \item No bloquea guardado (puede ser intencional).
        \item El usuario decide si continúa.
    \end{itemize}
\end{itemize}

\vspace{0.5em}
\noindent\textbf{Requerimientos especiales:}
\begin{itemize}[nosep]
    \item Cálculos automáticos en tiempo real.
    \item Validación de coherencia en pesos.
    \item Gestión dinámica de listados (agregar/editar/eliminar).
    \item Cambio de estado automático según rol.
    \item Registro automático en historial.
\end{itemize}

\subsubsection{Gestión de Análisis de PMS}

\noindent\textbf{Caso de uso:} Crear Repeticiones y Gestionar Tandas de PMS

\vspace{0.5em}
\noindent\textbf{Descripción:} El usuario crea repeticiones dentro de tandas, ingresa pesos, y el sistema calcula automáticamente promedios y coeficiente de variación, creando nuevas tandas si es necesario.

\vspace{0.5em}
\noindent\textbf{Pre-Condición:}
\begin{itemize}[nosep]
    \item El análisis de PMS debe existir con estado \texttt{EN\_PROCESO}
    \item La configuración inicial debe estar completa (repeticiones esperadas, tipo de semilla)
    \item Debe existir al menos una tanda
    \item El usuario debe estar autenticado
\end{itemize}

\vspace{0.5em}
\noindent\textbf{Post-Condición:}
\begin{itemize}[nosep]
    \item Repeticiones creadas con pesos registrados
    \item Cálculos automáticos realizados (promedio, CV)
    \item Si CV > umbral $\rightarrow$ nueva tanda creada automáticamente
    \item Análisis listo para finalizar cuando se cumplen condiciones
\end{itemize}

\vspace{0.5em}
\noindent\textbf{Actores:} Analista, Administrador.

\vspace{0.5em}
\noindent\textbf{Flujo de evento principal:}
\begin{enumerate}[nosep]
    \item \textbf{Fase 1: Acceso al análisis}
    \begin{itemize}[nosep]
        \item El usuario accede al detalle del análisis de PMS
        \item Ve la configuración establecida:
        \begin{itemize}[nosep]
            \item Número de repeticiones esperadas (4-20)
            \item Tipo de semilla (brozosa o normal)
            \item Umbral CV (6\% o 4\%)
        \end{itemize}
        \item Ve la tanda actual activa
    \end{itemize}
    \item \textbf{Fase 2: Creación de repeticiones}
    \begin{itemize}[nosep]
        \item El usuario selecciona ``Agregar repetición''
        \item Ingresa el peso en gramos (3 decimales)
        \item Guarda la repetición
        \item El sistema:
        \begin{itemize}[nosep]
            \item Crea la repetición
            \item La vincula a la tanda actual
            \item Incrementa el contador de repeticiones de la tanda
        \end{itemize}
    \end{itemize}
    \item \textbf{Fase 3: Cálculos automáticos}
    \begin{itemize}[nosep]
        \item Después de crear cada repetición, el sistema calcula:
        \begin{itemize}[nosep]
            \item Peso promedio de la tanda = suma de pesos / número de repeticiones
            \item Coeficiente de variación (CV) = (desviación estándar / promedio) $\times$ 100
        \end{itemize}
    \end{itemize}
    \item \textbf{Fase 4: Validación de CV y creación de tandas}
    \begin{itemize}[nosep]
        \item Si CV $\leq$ umbral:
        \begin{itemize}[nosep]
            \item La tanda queda válida
            \item Muestra indicador verde
        \end{itemize}
        \item Si CV > umbral:
        \begin{itemize}[nosep]
            \item El sistema crea automáticamente una nueva tanda
            \item Muestra mensaje: ``El CV excede el límite. Se creó automáticamente la Tanda [N]''
            \item La nueva tanda queda activa para agregar repeticiones
        \end{itemize}
        \item El proceso se repite hasta lograr CV $\leq$ umbral
    \end{itemize}
    \item \textbf{Fase 5: Finalización}
    \begin{itemize}[nosep]
        \item Cuando todas las repeticiones están completas y CV es válido:
        \begin{itemize}[nosep]
            \item El análisis queda listo para finalizar
            \item El usuario puede proceder a finalizar el análisis
        \end{itemize}
    \end{itemize}
\end{enumerate}

\vspace{0.5em}
\noindent\textbf{Flujos alternativos:}
\begin{itemize}[nosep]
    \item \textit{Intentar crear más repeticiones que las configuradas:}
    \begin{itemize}[nosep]
        \item El sistema valida el número de repeticiones esperadas
        \item Bloquea la creación si se excede
        \item Muestra mensaje: ``Ya se alcanzó el número máximo de repeticiones configuradas ([N])''
    \end{itemize}
    \item \textit{CV persistentemente alto:}
    \begin{itemize}[nosep]
        \item El sistema sigue creando tandas automáticamente
        \item El usuario puede revisar la configuración o datos ingresados
        \item Puede consultar con supervisor sobre la calidad de la muestra
    \end{itemize}
    \item \textit{Editar peso de repetición existente:}
    \begin{itemize}[nosep]
        \item El usuario puede editar repeticiones ya creadas
        \item El sistema recalcula automáticamente promedio y CV
        \item Puede generar nueva tanda si CV excede umbral
    \end{itemize}
    \item \textit{Eliminar repetición:}
    \begin{itemize}[nosep]
        \item El usuario puede eliminar repeticiones
        \item El sistema recalcula automáticamente
        \item Puede eliminar tandas vacías automáticamente
    \end{itemize}
\end{itemize}

\vspace{0.5em}
\noindent\textbf{Requerimientos especiales:}
\begin{itemize}[nosep]
    \item Cálculos automáticos de promedio y CV en tiempo real
    \item Creación automática de tandas cuando CV excede umbral
    \item No se puede modificar la configuración inicial (repeticiones esperadas, tipo de semilla)
    \item Validación de número máximo de repeticiones por tanda
    \item Pesos con 3 decimales de precisión
\end{itemize}

% 2.11 Caso de uso - Gestión de Análisis de Germinación
\subsubsection{Gestión de Análisis de Germinación}

\noindent\textbf{Caso de uso:} Crear y Gestionar Tablas de Germinación

\vspace{0.5em}
\noindent\textbf{Descripción:} El usuario crea tablas de germinación con condiciones experimentales específicas, agrega repeticiones, registra conteos por fecha, y el sistema calcula automáticamente totales y porcentajes.

\vspace{0.5em}
\noindent\textbf{Pre-Condición:}
\begin{itemize}[nosep]
    \item El análisis de Germinación debe existir con estado \texttt{EN\_PROCESO}
    \item El usuario debe estar autenticado
    \item El usuario debe tener permisos para crear tablas
\end{itemize}

\vspace{0.5em}
\noindent\textbf{Post-Condición:}
\begin{itemize}[nosep]
    \item Tabla(s) de germinación creada(s) con configuración específica
    \item Repeticiones creadas y conteos registrados
    \item Cálculos automáticos realizados (totales, promedios, porcentajes)
    \item Análisis listo para finalizar cuando hay al menos una tabla completa
\end{itemize}

\vspace{0.5em}
\noindent\textbf{Actores:} Administrador.

\vspace{0.5em}
\noindent\textbf{Flujo de evento principal:}
\begin{enumerate}[nosep]
    \item \textbf{Fase 1: Creación de tabla de germinación}
    \begin{enumerate}[nosep]
        \item El usuario accede al detalle del análisis de Germinación
        \item Selecciona ``Crear nueva tabla de germinación''
        \item Configura parámetros de la tabla:
        \begin{itemize}[nosep]
            \item Fecha de inicio
            \item Fechas de conteos (hasta 4 fechas)
            \item Número de repeticiones: entre 2 y 8
            \item Número de semillas por repetición: 25, 50 o 100
            \item Prefrío (días, opcional)
            \item Pretratamiento (descripción, opcional)
            \item Método: sustrato, entre papel, sobre papel, etc.
            \item Temperatura (°C)
            \item Tratamientos adicionales (opcional)
        \end{itemize}
        \item Guarda la configuración
        \item El sistema crea la tabla con estado inicial
    \end{enumerate}

    \item \textbf{Fase 2: Creación de repeticiones}
    \begin{enumerate}[nosep]
        \item El usuario selecciona ``Agregar repetición'' a la tabla
        \item El sistema crea la repetición (ej. Repetición 1, 2, 3...)
        \item La repetición queda lista para registrar conteos
        \item Repite hasta alcanzar el número de repeticiones configurado
    \end{enumerate}

    \item \textbf{Fase 3: Registro de conteos}
    \begin{enumerate}[nosep]
        \item Para cada repetición, el usuario ingresa conteos por fecha:
        \begin{itemize}[nosep]
            \item Fecha 1: número de semillas germinadas normalmente
            \item Fecha 2: número de semillas germinadas normalmente
            \item Fecha 3: número de semillas germinadas normalmente (si aplica)
            \item Fecha 4: número de semillas germinadas normalmente (si aplica)
        \end{itemize}
        \item Además puede registrar:
        \begin{itemize}[nosep]
            \item Semillas anormales
            \item Semillas muertas
            \item Semillas duras
            \item Semillas frescas
        \end{itemize}
    \end{enumerate}

    \item \textbf{Fase 4: Cálculos automáticos}
    \begin{enumerate}[nosep]
        \item El sistema calcula automáticamente por repetición y conteo:
        \begin{itemize}[nosep]
            \item Total de semillas germinadas por fecha
            \item Porcentaje de germinación = (germinadas / total semillas) $\\times$ 100
            \item Totales acumulados por conteo
        \end{itemize}
        \item Calcula promedios entre repeticiones:
        \begin{itemize}[nosep]
            \item Promedio de germinación normal por fecha
            \item Porcentaje promedio final
        \end{itemize}
    \end{enumerate}

    \item \textbf{Fase 5: Gestión de múltiples tablas}
    \begin{enumerate}[nosep]
        \item El usuario puede crear múltiples tablas con diferentes condiciones experimentales
        \item Cada tabla es independiente y tiene sus propios cálculos
        \item Permite comparar resultados entre diferentes condiciones
    \end{enumerate}
\end{enumerate}

\vspace{0.5em}
\noindent\textbf{Flujos alternativos:}
\begin{itemize}[nosep]
    \item \textit{Intentar finalizar análisis sin tablas:}
    \begin{itemize}[nosep]
        \item El sistema valida que exista al menos una tabla
        \item Bloquea finalización
        \item Muestra error: ``Debe crear al menos una tabla de germinación antes de finalizar''
    \end{itemize}

    \item \textit{Tabla sin repeticiones completas:}
    \begin{itemize}[nosep]
        \item El sistema no calcula promedios
        \item Muestra advertencia: ``Esta tabla tiene repeticiones incompletas''
        \item Permite continuar pero marca la tabla como incompleta
    \end{itemize}

    \item \textit{Conteo excede número de semillas por repetición:}
    \begin{itemize}[nosep]
        \item El sistema valida en tiempo real
        \item Muestra error: ``El conteo no puede exceder [N] semillas configuradas''
        \item Bloquea guardado hasta corregir
    \end{itemize}

    \item \textit{Editar configuración de tabla existente:}
    \begin{itemize}[nosep]
        \item El usuario puede editar parámetros de la tabla
        \item Si ya hay repeticiones con datos → muestra advertencia
        \        item Requiere confirmación para continuar
    \end{itemize}

    \item \textit{Eliminar tabla:}
    \begin{itemize}[nosep]
        \item El usuario puede eliminar tablas
        \item Si tiene repeticiones con datos → requiere confirmación
        \item Elimina en cascada todas las repeticiones y conteos asociados
    \end{itemize}
\end{itemize}

\vspace{0.5em}
\noindent\textbf{Requerimientos especiales:}
\begin{itemize}[nosep]
    \item Cada tabla es independiente con su propia configuración
    \item Cálculos automáticos en tiempo real por repetición y promediados
    \item Validación de conteos no excedan semillas configuradas
    \item Permite múltiples condiciones experimentales en un mismo análisis
    \item Fechas de conteo deben ser posteriores a fecha de inicio
    \item Número de fechas de conteo flexible (1 a 4)
\end{itemize}

% 2.12 Caso de uso - Gestión de Análisis de Tetrazolio
\subsection{Caso de uso: Gestión de Análisis de Tetrazolio}

\noindent\textbf{Caso de uso:} Crear Repeticiones y Registrar Conteos de Tetrazolio

\vspace{0.5em}
\noindent\textbf{Descripción:} El usuario crea repeticiones según la configuración establecida, registra conteos de semillas viables, no viables y muertas, y el sistema calcula automáticamente porcentajes de viabilidad.

\vspace{0.5em}
\noindent\textbf{Pre-Condición:}
\begin{itemize}[nosep]
    \item El análisis de Tetrazolio debe existir con estado \texttt{EN\_PROCESO}
    \item La configuración inicial debe estar completa
    \item El usuario debe estar autenticado
\end{itemize}

\vspace{0.5em}
\noindent\textbf{Post-Condición:}
\begin{itemize}[nosep]
    \item Repeticiones creadas con conteos registrados
    \item Cálculos automáticos realizados (totales, porcentajes, promedio general)
    \item Análisis listo para finalizar cuando todas las repeticiones están completas
\end{itemize}

\vspace{0.5em}
\noindent\textbf{Actores:} Analista, Administrador.

\vspace{0.5em}
\noindent\textbf{Flujo de evento principal:}
\begin{enumerate}[nosep]
    \item \textbf{Fase 1: Acceso al análisis}
    \begin{enumerate}[nosep]
        \item El usuario accede al detalle del análisis de Tetrazolio
        \item Ve la configuración establecida: fecha del ensayo, semillas por repetición (25, 50 o 100), número de repeticiones esperadas (2-8), pretratamiento, concentración, temperatura, tiempo, Viabilidad INASE (si se ingresó)
        \item Visualiza las repeticiones existentes (si las hubiera)
    \end{enumerate}

    \item \textbf{Fase 2: Creación de repeticiones}
    \begin{enumerate}[nosep]
        \item El usuario selecciona ``Agregar repetición''
        \item El sistema crea la repetición automáticamente con número secuencial (Rep. 1, 2, 3...) y número de semillas configurado
        \item La repetición queda lista para ingresar conteos
    \end{enumerate}

    \item \textbf{Fase 3: Registro de conteos}
    \begin{enumerate}[nosep]
        \item Para cada repetición, el usuario ingresa:
        \begin{itemize}[nosep]
            \item Semillas viables: cantidad de semillas con tinción adecuada
            \item Semillas no viables: cantidad de semillas con tinción inadecuada
            \item Semillas muertas: cantidad de semillas sin tinción
        \end{itemize}
        \item El sistema valida que la suma no exceda el número de semillas configurado
    \end{enumerate}

    \item \textbf{Fase 4: Cálculos automáticos}
    \begin{enumerate}[nosep]
        \item Por cada repetición, el sistema calcula:
        \begin{itemize}[nosep]
            \item Total evaluado = viables + no viables + muertas
            \item Porcentaje de viabilidad = (viables / total evaluado) $\\times$ 100
        \end{itemize}
        \item Calcula el promedio general de viabilidad entre todas las repeticiones
        \item Muestra comparación con Viabilidad INASE (si se ingresó)
    \end{enumerate}

    \item \textbf{Fase 5: Finalización}
    \begin{enumerate}[nosep]
        \item Cuando todas las repeticiones esperadas están completas, el análisis queda listo para finalizar
        \item Muestra resumen de resultados
        \item El usuario puede proceder a finalizar el análisis
    \end{enumerate}
\end{enumerate}

\vspace{0.5em}
\noindent\textbf{Flujos alternativos:}
\begin{itemize}[nosep]
    \item \textit{Conteo excede número de semillas:}
    \begin{itemize}[nosep]
        \item El sistema valida en tiempo real
        \item Muestra error: ``El total de conteos ([X]) excede las [N] semillas configuradas''
        \item Bloquea guardado hasta corregir
    \end{itemize}

    \item \textit{Intentar crear más repeticiones que las configuradas:}
    \begin{itemize}[nosep]
        \item El sistema valida el número de repeticiones esperadas
        \item Bloquea la creación
        \item Muestra mensaje: ``Ya se alcanzó el número máximo de repeticiones configuradas ([N])''
    \end{itemize}

    \item \textit{Editar conteos de repetición existente:}
    \begin{itemize}[nosep]
        \item El usuario puede editar repeticiones ya creadas
        \item El sistema recalcula automáticamente porcentajes y promedio
        \item Muestra cambios en tiempo real
    \end{itemize}

    \item \textit{Eliminar repetición:}
    \begin{itemize}[nosep]
        \item El usuario puede eliminar repeticiones
        \item El sistema recalcula automáticamente el promedio general
        \item Permite crear nuevas repeticiones hasta el límite configurado
    \end{itemize}

    \item \textit{Conteo incompleto:}
    \begin{itemize}[nosep]
        \item El usuario intenta guardar sin completar todos los campos
        \item El sistema muestra advertencia
        \item Permite guardar como borrador pero marca repetición como incompleta
    \end{itemize}
\end{itemize}

\vspace{0.5em}
\noindent\textbf{Requerimientos especiales:}
\begin{itemize}[nosep]
    \item No se puede modificar la configuración inicial una vez creadas las repeticiones
    \item Validación estricta de suma de conteos $\leq$ semillas configuradas
    \item Cálculos automáticos de porcentajes con 2 decimales
    \item Comparación automática con Viabilidad INASE
    \item Todas las repeticiones deben estar completas antes de finalizar
\end{itemize}

% 2.13 Caso de uso - Gestión de Análisis DOSN
\subsubsection{Caso de uso: Gestión de Análisis DOSN}

\noindent\textbf{Caso de uso:} Editar y Completar Análisis DOSN

\vspace{0.5em}
\noindent\textbf{Descripción:} El usuario edita un análisis DOSN existente, completa o modifica datos INIA e INASE, gestiona listados de malezas, otros cultivos, brassica y cuscuta, y define cumplimiento del estándar.

\vspace{0.5em}
\noindent\textbf{Pre-Condición:}
\begin{itemize}[nosep]
    \item El análisis DOSN debe existir
    \item El análisis debe estar en estado \texttt{EN\_PROCESO} o \texttt{PENDIENTE\_APROBACION}
    \item Si está \texttt{APROBADO}, solo Administrador puede editar
    \item El usuario debe estar autenticado
\end{itemize}

\vspace{0.5em}
\noindent\textbf{Post-Condición:}
\begin{itemize}[nosep]
    \item Datos del análisis actualizados
    \item Si Analista edita \texttt{APROBADO} → cambia a \texttt{PENDIENTE\_APROBACION}
    \item Si Administrador edita \texttt{APROBADO} → permanece \texttt{APROBADO}
    \item Registro automático en historial de modificación
\end{itemize}

\vspace{0.5em}
\noindent\textbf{Actores:} Analista, Administrador.

\vspace{0.5em}
\noindent\textbf{Flujo de evento principal:}
\begin{enumerate}[nosep]
    \item \textbf{Acceso al análisis}
    \begin{enumerate}[nosep]
        \item El usuario accede al detalle del análisis DOSN
        \item Selecciona ``Editar análisis''
        \item El sistema carga los datos actuales en dos pestañas:
        \begin{itemize}[nosep]
            \item Datos generales (INIA e INASE)
            \item Registros (listados)
        \end{itemize}
    \end{enumerate}

    \item \textbf{Edición de datos INIA}
    \begin{enumerate}[nosep]
        \item El usuario puede modificar: fecha de análisis, gramos analizados, tipos de análisis (Completo, Reducido, Limitado, Reducido-Limitado)
        \item El sistema valida: fecha no futura, gramos > 0, al menos un tipo de análisis marcado
    \end{enumerate}

    \item \textbf{Edición de datos INASE (opcionales)}
    \begin{enumerate}[nosep]
        \item El usuario puede modificar: fecha INASE, gramos analizados INASE, tipos de análisis INASE
        \item Si no completa, los datos INASE permanecen vacíos (opcionales)
    \end{enumerate}

    \item \textbf{Gestión de listados de Malezas}
    \begin{enumerate}[nosep]
        \item El usuario puede agregar nuevas malezas identificadas
        \item Seleccionar tipo: Tolerancia cero, Toleradas, o Comunes
        \item Seleccionar catálogo o especie e ingresar cantidad de semillas
        \item Editar o eliminar registros existentes
    \end{enumerate}

    \item \textbf{Gestión de Otros Cultivos}
    \begin{enumerate}[nosep]
        \item Agregar otros cultivos identificados, seleccionar catálogo o especie e ingresar cantidad
        \item Editar o eliminar registros existentes
    \end{enumerate}

    \item \textbf{Gestión de Brassica}
    \begin{enumerate}[nosep]
        \item Si aplica, agregar semillas brassica identificadas, seleccionar catálogo o especie e ingresar cantidad
        \item Editar o eliminar registros existentes
    \end{enumerate}

    \item \textbf{Registro de Cuscuta}
    \begin{enumerate}[nosep]
        \item Si se encuentra cuscuta, el usuario registra cantidad de semillas de cuscuta y observaciones
    \end{enumerate}

    \item \textbf{Cumplimiento del estándar}
    \begin{enumerate}[nosep]
        \item El usuario selecciona cumplimiento (Sí/No) y guarda los cambios
        \        item El sistema actualiza el análisis, registra la modificación en historial y aplica cambio de estado según rol
    \end{enumerate}
\end{enumerate}

\vspace{0.5em}
\noindent\textbf{Flujos alternativos:}
\begin{itemize}[nosep]
    \item \textit{Analista edita análisis APROBADO}
    \begin{itemize}[nosep]
        \item El sistema detecta que es Analista
        \item Permite editar
        \item Al guardar, cambia estado a ``PENDIENTE APROBACIÓN''
        \item Muestra mensaje de cambio de estado
    \end{itemize}

    \item \textit{Administrador edita análisis APROBADO}
    \begin{itemize}[nosep]
        \item El sistema detecta que es Administrador
        \item Permite editar
        \item Al guardar, mantiene estado ``APROBADO''
    \end{itemize}

    \item \textit{Fecha futura en INIA}
    \begin{itemize}[nosep]
        \item El sistema valida fecha INIA
        \item Muestra error: ``Ingrese una fecha válida (no futura)''
        \item No permite guardar hasta corregir
    \end{itemize}

    \item \textit{Gramos inválidos}
    \begin{itemize}[nosep]
        \item El sistema valida que gramos INIA > 0
        \item Muestra error correspondiente
        \item No permite guardar hasta corregir
    \end{itemize}

    \item \textit{Sin tipo de análisis INIA}
    \begin{itemize}[nosep]
        \item El sistema valida al menos un checkbox marcado
        \item Muestra error: ``Debe seleccionar al menos un tipo de análisis''
        \item No permite guardar hasta corregir
    \end{itemize}

    \item \textit{Eliminar todos los listados}
    \begin{itemize}[nosep]
        \item El usuario puede eliminar todos los registros
        \item El análisis queda solo con datos INIA/INASE
        \item El sistema permite pero advierte antes de finalizar
    \end{itemize}
\end{itemize}

\vspace{0.5em}
\noindent\textbf{Requerimientos especiales:}
\begin{itemize}[nosep]
    \item Datos INIA obligatorios, datos INASE opcionales
    \item Validación de fechas no futuras
    \item Gestión dinámica de múltiples listados independientes
    \item Cada tipo de listado (malezas, cultivos, brassica) es independiente
    \item Cambio de estado automático según rol
    \item Registro automático en historial
\end{itemize}

